\section{Models and Data}
\label{sec:models-and-data}

\subsection{Models}
\label{sec:models}

We study five AI weather models from three architecture families
(Table~\ref{tab:models}). They were chosen
from the models whose trained weights are publicly available and that
can be run in PyTorch on a single GPU with 16\,GB of memory, to cover
the main architecture families in current use: transformers (Aurora, Pangu-Weather), models that process
the fields in spectral space (FourCastNet, SFNO), and a graph neural
network (FastNet). All are
released with fp32 weights and run in fp32 by default; the fp32 forecast
is therefore the reference against which every reduced-precision
forecast is compared. All except FastNet run on the regular
0.25$^\circ$ latitude--longitude grid of ERA5, and all except
Pangu-Weather advance the forecast in 6-hour steps.

\begin{table}[htbp]
\centering
\footnotesize
\caption{Models studied. Variables: number of surface variables plus
number of atmospheric variables $\times$ number of pressure levels
(FourCastNet: 18 atmospheric fields on selected levels, see text; Aurora:
inputs, with seven further surface variables output only).
Runtime: time per forecast step in fp32 on our 16\,GB GPU
(Table~\ref{tab:benchmark}).}
\label{tab:models}
\setlength{\tabcolsep}{4pt}
\begin{tabular}{lllllrrl}
\toprule
Model & Version & Architecture & Grid & Step & Variables & Params & Runtime\\
\midrule
\shortstack[l]{FastNet\\\citep{dunstan2025fastnet}} & v1.1 & graph network & O96 ($\sim$1$^\circ$) & 6h & $7 + 5{\times}13$ & 62.2M & 0.29\,s \\
\shortstack[l]{FourCastNet\\\citep{pathak2022fourcastnet}} & v1 & Fourier operator & 0.25$^\circ$ & 6h & $8 + 18$ & 75.3M & 0.12\,s \\
\shortstack[l]{Pangu-Weather\\\citep{bi2023pangu}} & 24h model & 3D transformer & 0.25$^\circ$ & 24h & $4 + 5{\times}13$ & 276.7M & 0.79\,s \\
\shortstack[l]{SFNO\\\citep{bonev2023sfno}} & v0.1.0 & spherical Fourier op. & 0.25$^\circ$ & 6h & $8 + 5{\times}13$ & 289.4M & 0.53\,s \\
\shortstack[l]{Aurora\\\citep{bodnar2024aurora}} & 1.5 det./ens. & 3D Swin transformer & 0.25$^\circ$ & 6h & $18 + 5{\times}13$ & 1.26B & 5.4\,s$^\dagger$ \\
\bottomrule
\end{tabular}\\[2pt]
{\scriptsize $^\dagger$~With the backbone
weights streamed through GPU memory, the only way Aurora fits in 16\,GB
in fp32 (\S\ref{sec:quantisation}, \S\ref{sec:fp32-cost}).}
\end{table}

Aurora \citep{bodnar2024aurora} is a 1.26-billion-parameter model
pre-trained on a large and diverse set of weather and climate data and
then fine-tuned for forecasting. An encoder maps the input fields -- two
consecutive analysis times, 18 surface variables and five atmospheric
variables on 13 pressure levels -- to a compact internal representation;
a backbone of 3D Swin-transformer layers, which compute interactions
within local windows that shift from layer to layer, advances this
representation by one 6-hour step; and a decoder maps it back to the
physical variables. We use version~1.5 \citep{aurora_docs} in two
checkpoints from the same release: a deterministic one and an ensemble
one. The two have almost the same architecture and training data, but
the ensemble checkpoint adds random noise inside the model, so that
different random seeds give different ensemble members, and was fine-tuned
with the continuous ranked probability score (CRPS) as its loss instead
of the mean absolute error (MAE) \citep{weyn2026aurora15}.

Pangu-Weather \citep{bi2023pangu} is a 3D transformer with a
positional encoding adapted to the Earth's geometry, operating on four
surface variables and five atmospheric variables on 13 pressure levels.
Separate models were trained for 1-, 3-, 6- and 24-hour steps; we use
the 24-hour model, so its forecasts advance in 24-hour steps. The
official release provides the model only as exported ONNX computation
graphs, not as PyTorch code that can be converted to 16 bits component by
component, so we use a third-party PyTorch reimplementation of the
24-hour model with the official weights \citep{pangu_pytorch}.

FourCastNet \citep{pathak2022fourcastnet} (version~1) has 26 variables in the released
checkpoint (the original paper describes 20):
8 surface or single-level fields (10\,m and 100\,m wind, 2\,m temperature,
surface and mean sea-level pressure, total column water vapour) and 18
atmospheric fields, with each variable on its own subset of levels:
geopotential at 1000, 850, 500, 250 and 50\,hPa; the wind components at
1000, 850, 500 and 250\,hPa; temperature at 850, 500 and 250\,hPa; and
relative humidity at 850 and 500\,hPa. It
mixes information between grid points with an Adaptive Fourier Neural
Operator \citep{guibas2022afno}, which multiplies the Fourier transform
of the internal fields by learned complex-valued filters. The Fourier
transform treats the latitude--longitude grid as a flat plane rather
than a sphere. Although the filters are complex-valued, the
implementation stores their real and imaginary parts as separate
real-valued weights and carries out the complex multiplication with
real arithmetic, so the model can be converted to 16 bits, with one
small change to its code.

FastNet \citep{dunstan2025fastnet, daub2025fastnet}, developed by the
Met Office and the Alan Turing Institute (version~1.1), is a graph
neural network of the GraphCast family \citep{lam2023graphcast}. Its
encoder passes information from the grid to an icosahedral mesh of
10{,}242 nodes in a single message-passing step; its processor performs
16 rounds of message passing on that mesh, each with its own weights;
and its decoder passes the result back to the grid, adding it to the
input state. It forecasts seven surface variables and five atmospheric
variables on 13 pressure levels, with 12 further static and time-of-day
fields as input. Unlike the other models, it runs on ECMWF's O96 reduced
Gaussian grid (40{,}320 points, roughly 1$^\circ$ resolution). We build
its initial conditions from the same 0.25$^\circ$ ERA5 data as for the
other models, by linear interpolation onto the O96 points, and
interpolate its forecasts back to the 0.25$^\circ$ grid before
verification (\S\ref{sec:verification-data}), so that all models are
scored on the same grid against the same truth.

SFNO \citep{bonev2023sfno} (the \texttt{sfno\_73ch\_small}
checkpoint, version~0.1.0) replaces FourCastNet's flat Fourier transform
with a transform in spherical harmonics, which respects the geometry of
the sphere. It forecasts eight surface variables and five atmospheric
variables on 13 pressure levels. Its spectral filters, 97.8\% of its
parameters, are complex-valued; in 16 bits we apply them as real and
imaginary parts, as FourCastNet does.

We also considered two further models. ECMWF's AIFS \citep{lang2024aifs}
(\texttt{aifs-single-2.0}) runs on our GPU, but it needs six ocean-wave
variables as input that ERA5 does not archive, so it cannot be
initialised like the other models. GraphCast \citep{lam2023graphcast} is
implemented in JAX, and no PyTorch version loads its official weights.

Like most deep-learning models, AI weather models can be divided into
components. In a forecast step, the physical fields of the state are
typically normalised, passed through an encoder, a number of backbone
layers or blocks, and a decoder, and the decoder's output is
denormalised into the physical fields of the next state. We use this
division in the experiments that locate where fp16 rounding enters the
forecast (\S\ref{sec:methods-tier2}). The components differ widely in
their numbers of parameters (Table~\ref{tab:components}).

\begin{table}[htbp]
\centering
\small
\caption{Components per model, in forward-pass order, with parameter counts and share of the model total. Repeated identical blocks are grouped. ``Other'': parameters outside every component, which stay fp32 in the one-at-a-time experiments (\S\ref{sec:methods-tier2}) (FourCastNet: its learned positional embedding; Aurora: backbone time/noise MLPs and context-downsampling layers; SFNO: the $1\times1$ convolution on its residual path). SFNO's complex filter weights are counted once each.}
\label{tab:components}
\begin{tabular}{llrr}
\toprule
Model & Component & Parameters & Share \\
\midrule
Aurora & \texttt{encoder} & 7{,}665{,}152 & 0.608\% \\
 & \texttt{backbone\_stage\_0} & 27{,}307{,}008 & 2.17\% \\
 & \texttt{backbone\_stage\_1} & 155{,}330{,}560 & 12.3\% \\
 & \texttt{backbone\_stage\_2} & 436{,}420{,}608 & 34.6\% \\
 & \texttt{backbone\_stage\_3} & 445{,}859{,}840 & 35.4\% \\
 & \texttt{backbone\_stage\_4} & 149{,}294{,}080 & 11.8\% \\
 & \texttt{backbone\_stage\_5} & 25{,}205{,}760 & 2\% \\
 & \texttt{decoder} & 9{,}953{,}776 & 0.79\% \\
 & other & 3{,}147{,}776 & 0.25\% \\
 & \textit{total} & 1{,}260{,}184{,}560 & \\
\addlinespace
Pangu-Weather & \texttt{patch\_embedding} & 58{,}752 & 0.0212\% \\
 & \texttt{stage\_0} & 31{,}744{,}896 & 11.5\% \\
 & \texttt{downsample} & 296{,}448 & 0.107\% \\
 & \texttt{stage\_1..2 (2$\times$)} & 106{,}198{,}272 & 38.4\% each \\
 & \texttt{upsample} & 332{,}160 & 0.12\% \\
 & \texttt{stage\_3} & 31{,}744{,}896 & 11.5\% \\
 & \texttt{patch\_recovery} & 86{,}240 & 0.0312\% \\
 & \textit{total} & 276{,}659{,}936 & \\
\addlinespace
FourCastNet & \texttt{patch\_embed} & 1{,}278{,}720 & 1.7\% \\
 & \texttt{block\_0..11 (12$\times$)} & 5{,}023{,}488 & 6.67\% each \\
 & \texttt{head} & 1{,}277{,}952 & 1.7\% \\
 & other & 12{,}441{,}600 & 16.5\% \\
 & \textit{total} & 75{,}280{,}128 & \\
\addlinespace
SFNO & \texttt{encoder} & 177{,}408 & 0.0613\% \\
 & \texttt{block\_0..7 (8$\times$)} & 36{,}129{,}408 & 12.5\% each \\
 & \texttt{decoder} & 175{,}872 & 0.0608\% \\
 & other & 5{,}621 & 0.0019\% \\
 & \textit{total} & 289{,}394{,}165 & \\
\addlinespace
FastNet & \texttt{mesh\_edge\_embedder} & 594{,}432 & 0.955\% \\
 & \texttt{encoder} & 1{,}843{,}200 & 2.96\% \\
 & \texttt{processor\_0..15 (16$\times$)} & 3{,}545{,}088 & 5.7\% each \\
 & \texttt{decoder} & 3{,}077{,}448 & 4.94\% \\
 & \textit{total} & 62{,}236{,}488 & \\
\bottomrule
\end{tabular}
\end{table}


In fp32, Aurora's deterministic checkpoint is the most accurate model
at 1 day, Pangu-Weather is level with it by 10 days, and FourCastNet and
SFNO are the least accurate from day 5 on (Figure~\ref{fig:model-skill};
all models verified on the same 52 initialisations, ERA5 truth and
0.25$^\circ$ grid, \S\ref{sec:verification-data}). At 1 day, Aurora has
the lowest global RMSE for every variable, with Pangu-Weather second; at
10 days, Pangu-Weather is level with or marginally better than Aurora for
every variable. FourCastNet and SFNO have the highest RMSE at days 5 and
10 for every variable; FourCastNet's 2t error grows most by day 10, where
a few of its fp32 forecasts blow up (\S\ref{sec:results-fp16}). FastNet
is among the least skilful models at day 1 (last for 2t and 10u) but
second or third by day 10. This is plausibly a resolution effect: its native grid
is about 1$^\circ$, and interpolation to 0.25$^\circ$ cannot restore
small-scale detail, which matters most at short lead times. Aurora's
ensemble checkpoint is verified as an ensemble in
\S\ref{sec:results-ensemble}. Figure~\ref{fig:model-skill-uk}
(Appendix~\ref{sec:appendix-figures}) shows the same comparison for the
UK/Europe region.

\begin{figure}[htbp]
\centering
\includegraphics[width=\textwidth]{model_comparison_rmse_global}
\caption{fp32 forecast skill of the five models (Aurora: deterministic checkpoint): mean global RMSE against ERA5 at daily lead times, 52 initialisations in 2024, all models verified on the regular 0.25$^\circ$ grid (FastNet's forecasts linearly interpolated from its O96 grid). Shading: $\pm$1 standard error across dates. UK/Europe: Figure~\ref{fig:model-skill-uk} in
Appendix~\ref{sec:appendix-figures}.}
\label{fig:model-skill}
\end{figure}

\label{sec:fp32-cost}
In fp32, a forecast step takes 0.12\,s (FourCastNet), 0.29\,s
(FastNet), 0.53\,s (SFNO), 0.79\,s (Pangu-Weather, a 24-hour step) and
5.4\,s (Aurora) on our GPU (Table~\ref{tab:benchmark},
\S\ref{sec:methods-further}), so that a 10-day forecast takes between 5\,s
and 3.6\,minutes and uses between 0.4 and 18\,Wh of energy. Peak memory
ranges from 1.2\,GB (FourCastNet) to 15.4\,GB (SFNO), which leaves SFNO
little headroom on a 16\,GB GPU. Aurora does not fit in 16\,GB in fp32 at
all: its fp32 runs move the backbone's weights to the GPU one block at a
time (\S\ref{sec:quantisation}), which keeps its peak memory at 14.3\,GB
but makes it the slowest model by far. These costs are the reference for
the reduced-precision runs (\S\ref{sec:results-fp16}).

\subsection{Quantisation}
\label{sec:quantisation}

A floating-point number is stored as a sign, an exponent and a
significand (mantissa). fp32 uses 8 exponent and 23 mantissa bits, so
that rounding a value to fp32 changes it by at most about
$6\times10^{-8}$ of its magnitude. The two common 16-bit formats divide
their bits differently. fp16 uses 5 exponent and 10 mantissa bits:
rounding changes a value by at most $2^{-11}\approx0.05\%$, but values
above 65{,}504 cannot be represented and become infinite, and values below
$6.1\times10^{-5}$ lose precision. bf16 keeps fp32's 8 exponent bits, and
hence its range (about $10^{-38}$ to $3\times10^{38}$), but only 7
mantissa bits, so that rounding changes a value by up to
$2^{-8}\approx0.4\%$, eight times more than fp16. bf16 can therefore
replace fp32 without regard to the size of the values
\citep{kalamkar2019bfloat16}; fp16 is more precise, but only if the values
stay within its range.

The simplest way to run a model in 16 bits is to convert the whole
model with one command (in PyTorch, \texttt{model.to(torch.bfloat16)} or
\texttt{model.to(torch.float16)}). This converts the weights, but a layer
can only multiply inputs that have the same precision as its weights, so
the inputs must be converted too. In models that normalise their inputs
themselves, this happens before the normalisation: Aurora's code
converts the input fields to the precision of its weights and then
normalises them, and Pangu-Weather normalises inside its first layer,
which therefore receives 16-bit inputs. The physical fields themselves
are then rounded to 16 bits at every step, which in fp16 overflows and in
bf16 costs substantial skill (\S\ref{sec:results-state}). In all our
16-bit runs, the state, its normalisation and denormalisation, including
the normalisation statistics, therefore stay in fp32, and only the
network is converted.

In the \emph{fp16} and \emph{bf16} runs, every weight and every
computation of the network is in that format, while the normalisation of
the inputs, the denormalisation of the outputs and the state carried from
one forecast step to the next stay in fp32. The normalised fields are rounded to 16
bits when they enter the network, and its output is converted back to
fp32 before it is denormalised. For Aurora, the input fields are kept and
normalised in fp32 and rounded on entry to the encoder; for
Pangu-Weather, the normalisation is moved out of the first layer; FastNet
carries a normalised state, which is kept in fp32 between steps;
FourCastNet's code already normalises and denormalises outside the
network in fp32; and for SFNO, the solar zenith angle and static fields
are added to the state in fp32 before it enters the network. We verified that the state never reaches the
network in 16 bits, and, for Pangu-Weather, that the modified code with
fp32 weights reproduces the original model bit for bit. Every weight is
rounded once -- in bf16, Aurora's weights change by 0.14\% on average --
and every intermediate result inside the network is stored in 16 bits.
Because current GPUs sum the products in a 16-bit matrix multiplication
in fp32, the main source of error within the network is the rounding of
these stored values rather than the arithmetic within an operation. A
few operations are computed in fp32 on 16-bit inputs: FourCastNet's
Fourier transforms and SFNO's spherical harmonic transforms. FourCastNet's network is converted component by
component (Table~\ref{tab:components}), which leaves its positional
embedding, part of no component, in fp32. Further numerical details are
given in the code's \texttt{nwpprec/README.md} \citep{siegert2026code}.

The fp32 reference is each model as PyTorch runs it by default. This
includes one deviation from strict fp32 that we kept deliberately:
PyTorch computes convolutions, which occur in the input and output
layers of Aurora, Pangu-Weather and FourCastNet, in the reduced-precision
format TF32 by default. Aurora
does not fit in 16\,GB of GPU memory in fp32; its fp32 runs therefore
move the backbone weights to the GPU one block at a time and compute
the decoder in spatial portions. This makes them slower
(\S\ref{sec:fp32-cost}) but, as we verified on test data, gives
bit-identical results.

The approach used here is a \emph{persistent} conversion: the model is
converted once and the fp32 weights are discarded. The most common
alternative, automatic mixed precision \citep{micikevicius2018mixed}, is the
default inference mode of some models, such as Aurora~1.5. It keeps the
weights in fp32 and converts values to 16 bits only for individual operations.
This protects accuracy but saves little memory: in this mode, Aurora~1.5 runs
out of memory on our 16\,GB GPU in the first forecast step, whereas the
persistently converted model needs at most 11.6\,GB
(Table~\ref{tab:benchmark}). The persistent conversion is the harsher of the
two for accuracy, so the costs reported here are those of the configuration
that fits on such a GPU.

A larger literature concerns quantisation to integers (int8, int4),
in which values are represented as integers multiplied by a scale shared
by a tensor, a channel or a group of values
\citep{jacob2018quantization, nagel2021white}. The scales are chosen from calibration data
after training (post-training quantisation) or learned during training
(quantisation-aware training), and either the weights alone or also the
intermediate results are quantised. \citet{bhattacharya2026ptq} applied
such methods, including ones developed for language models
\citep{xiao2022smoothquant, lin2023awq}, to two AI weather models.
8-bit floating-point formats (fp8) also need scaling
\citep{micikevicius2022fp8}. We did not pursue these formats: they need
calibration and low-precision implementations of every operation, and in
tests with Aurora they were either slower than bf16 or did not fit in
16\,GB. fp16 and bf16 are the
reduced-precision formats that most users of these models can apply
directly: they need no calibration, and fast fp16 arithmetic is
available on NVIDIA GPUs from 2017 on, bf16 arithmetic from 2020 on
\citep{nvidia2017volta, nvidia2020ampere}, and both on current consumer and
data-centre models. On CPUs, by contrast, fast fp16 arithmetic is
not usually available, so the speedups reported here should not be
expected to carry over to them.


\subsection{Model verification}
\label{sec:verification-data}

All forecasts are initialised from, and verified against, the ERA5
reanalysis \citep{hersbach2020era5} on its regular 0.25$^\circ$
latitude--longitude grid ($721\times1440$ points), the reference
commonly used to evaluate AI weather models
\citep{rasp2020weatherbench}. Models that take two input times (Aurora)
receive the analyses at the initial time and six hours earlier.
The grids of Aurora and FourCastNet omit the southernmost row of the ERA5
grid, which is therefore not scored for these models. Aurora~1.5 was fine-tuned last on
operational IFS analyses rather than ERA5 \citep{aurora_docs}; starting
it from ERA5 may lower its fp32 skill slightly, but affects its fp32
and fp16 forecasts alike.

Initial dates are drawn from 2024, a year after the training periods of
the models for which these are documented (\S\ref{sec:methods-tier1}
describes the sampling). 2024 lies after the training periods of
Pangu-Weather (1979--2017; \citealp{bi2023pangu}), FourCastNet
(1979--2015; \citealp{pathak2022fourcastnet}) and FastNet (1980--2021;
\citealp{daub2025fastnet}), and was chosen over 2023 because Aurora's
training data include 2023: Aurora~1.5 was trained on ERA5 from
1981--2023 and fine-tuned on operational analyses from 2018--2023, with
2024 held out as its test year \citep{weyn2026aurora15}. The SFNO
checkpoint's documentation gives no training period; the SFNO paper,
which it cites, trains its 73-variable models on 1979--2015 with 2016 and
2017 for validation \citep{bonev2023sfno}, the same split into training,
validation and out-of-sample data that the checkpoint's stored
configuration records, and the developer's SFNO inference service lists ERA5
from 1979--2017. Either way, 2024 lies outside it.

FastNet runs on the O96 grid of about 1$^\circ$ resolution
(\S\ref{sec:models}). Its initial conditions are obtained by linear
interpolation of the 0.25$^\circ$ ERA5 fields onto the O96 points, and
its forecasts are interpolated back to the 0.25$^\circ$ grid by linear
interpolation on the Delaunay triangulation of the O96 points, so that
all models are scored on the same grid against the same truth. The grid
points outside the triangulation, about 1.6\% (a gap at the 0$^\circ$
meridian and the rows at the poles), have no interpolated value and are not
scored for FastNet. Every
FastNet score in this paper goes through this step.

We verify 2\,m temperature (2t), mean sea-level pressure (msl),
10\,m wind speed (ws10, computed from the two wind components),
geopotential at 500\,hPa (z500) and temperature at 850\,hPa (t850), at
daily lead times from 1 to 10 days; Figure~\ref{fig:model-skill} shows
the two wind components instead of the wind speed. Since the lead times
are multiples of 24 hours, all models, including Pangu-Weather with its
24-hour steps, are verified at the same times. Each forecast is scored
globally, over three regions and at five locations
(\S\ref{sec:methods-regions}).

The main score is the root-mean-square error (RMSE) of each forecast
against ERA5, with each grid point weighted by the cosine of its
latitude to account for the convergence of the meridians
\citep{rasp2020weatherbench}. We also record the bias and the mean
absolute error (MAE), the deterministic counterpart of the CRPS. Scores
are computed per forecast and averaged over initial dates. The ensemble
forecasts (\S\ref{sec:results-ensemble}) are additionally scored by the
area-weighted mean of the fair continuous ranked probability score
(CRPS), which does not penalise ensembles for having few members
\citep{ferro2014fair}; by the ensemble spread, the square root of the
area mean of the ensemble variance, compared with the RMSE of the
ensemble mean after correcting for the finite ensemble size
\citep{leutbecher2008ensemble}; and by the fraction of grid points at
which ERA5 lies outside the range of the ensemble, which for a
calibrated ensemble of $M$ members is $2/(M+1)$ \citep{siegert2011outliers}.
